% ============================================================
% CAPITOLO 9 - CONCLUSIONI E RINGRAZIAMENTI
% ============================================================

\section{Conclusioni e Sviluppi Futuri}

Il progetto \textbf{Signify} ha dimostrato la fattibilità di un sistema di riconoscimento della lingua dei segni americana (ASL) basato esclusivamente su computer vision e architetture neurali leggere, eliminando la necessità di sensori hardware costosi o ingombranti. L'integrazione di \textit{MediaPipe Holistic} per l'estrazione dei landmark e di una rete \textit{LSTM con meccanismo di Attention} ha permesso di ottenere prestazioni solide nella classificazione di segni isolati, gettando le basi per strumenti di comunicazione realmente inclusivi.

L'analisi dei risultati ha confermato che, mentre il solo modello visivo può presentare ambiguità tra gesti simili, l'approccio ibrido con correzione semantica (NLP) è la chiave per elevare l'accuratezza a livelli utilizzabili nel quotidiano.

\subsection{Sviluppi Futuri}
Nonostante i traguardi raggiunti, il percorso verso una traduzione fluida e universale presenta ancora ampi margini di miglioramento:

\begin{itemize}
    \item \textbf{Ottimizzazione dell'estrazione Live-Web:} Attualmente, l'estrazione dei punti chiave tramite \textit{Mediapipe} in ambiente browser può risultare carente o instabile in condizioni di scarsa illuminazione o con background complessi. Future iterazioni potrebbero prevedere l'uso di modelli di estrazione personalizzati e ottimizzati tramite WebAssembly (WASM) per una maggiore stabilità temporale dei landmark.
    \item \textbf{Modelli NLP avanzati per la frase:} L'integrazione NLP presentata come \textit{Proof of Concept} può essere estesa a veri e propri modelli di \textit{Neural Machine Translation} (NMT) capaci di gestire non solo la correzione lessicale, ma anche la complessa grammatica della lingua dei segni (che differisce sensibilmente da quella parlata), permettendo la composizione di frasi articolate e naturali.
\end{itemize}

\vspace{2cm}

\section*{Ringraziamenti}
\addcontentsline{toc}{section}{Ringraziamenti}

Giunto alla conclusione di questo percorso, desidero rivolgere un sincero ringraziamento ai professori del corso di \textbf{Machine Learning} della laurea triennale in Informatica dell'Università degli Studi di Salerno. Gli insegnamenti e le competenze trasmesse durante il corso sono stati fondamentali per fornirmi gli strumenti tecnici e teorici necessari per affrontare e modellare la complessità di questo progetto.

Un ringraziamento speciale va a chiunque abbia dedicato del tempo per leggere questo lavoro. Questo progetto nasce dalla volontà di unire tecnologia e inclusività, con l'obiettivo di fornire un piccolo contributo alla risoluzione di una problematica reale che coinvolge milioni di persone nel mondo.
